% from ega4-15-fr.tex
\oldpage[IV-3]{242}
Considérons d'autre part le point fermé $x_2$ de $X_2$ correspondant à l'idéal maximal $\left(u-\frac{1}{\pi}\right)+(v)$ de $K[u,v]$~; on a aussi $\mathbf{k}(x_2)=K$. Comme on le verra au chap.~V, on peut donc \emph{recoller} $X_1$ et $X_2$ le long des deux sous-préschémas fermés $Z_1$, $Z_2$ de $X_1$, $X_2$ ayant respectivement $\{x_1\}$ et $\{x_2\}$ pour espaces sous-jacents, suivant l'unique $Y$-isomorphisme de $Z_1$ sur $Z_2$. Désignons par $X$ le $Y$-préschéma ainsi obtenu, par $x$ le point de $X$ image de $x_1$ et $x_2$. La «~section nulle~» $g$ de $X_1$ (\textbf{II}, 1.7.9) est encore alors une $Y$-section de $X$~; $f^{-1}(y)$ est égale à $(X_1)_y$, tandis que $f^{-1}(y')$ a deux composantes irréductibles, respectivement isomorphes à $(X_1)_{y'}^0$ et $(X_2)_{y'}^0$, ayant un unique point commun $x$, et de dimensions respectives 1 et 2. Voir cependant la prop.~(15.6.9).

(ii) Considérons le morphisme $f$ défini dans (12.2.3, (ii)), qui est propre et plat~; en outre, la restriction de $f$ à $X_1\cap f_0^{-1}(Y)$ est un isomorphisme, donc le morphisme réciproque $g:Y\to X$ est une $Y$-section de $X$. On a alors $\dim(X_{y_0})=1$ tandis que $\dim(X_y^0)=0$ pour $y\ne y_0$, bien que $X_y^0$ soit géométriquement irréductible pour tout $y\in Y$ (mais $X_{y_0}^0$ n'est pas réduit)~; on voit donc que dans (15.6.6, (iii)), on ne peut supprimer les hypothèses $\alpha$) et $\beta$). En outre, $X^0$ n'est pas un voisinage de $X_{y_0}^{0}$, ce qui prouve que dans (15.6.4), on ne peut se dispenser de l'hypothèse que $X_{y_0}^0$ est réduit.

(iii) Au chap.~VI, nous appliquerons les résultats précédents aux $Y$-préschémas \emph{en groupes} localement de type fini sur un préschéma localement noethérien $Y$. Si $G$ est un tel préschéma, il existe une $Y$-section canonique $g$, la «~section unité~», et l'on montre que $G_y^0$ est toujours géométriquement irréductible et que l'on a $\dim(G_y^0)=\dim(G_y)$, en posant $G_y=f^{-1}(y)$. Il résultera donc de (15.6.2) et (15.6.3) que la fonction $z\mapsto\dim(G_z)$ est semi-continue supérieurement, et que si elle est constante au voisinage d'un point $y$, $f$ est universellement ouvert aux points de $G_y^0$. Si de plus $G_y$ est géométriquement réduit sur $\mathbf{k}(y)$ en un de ses points, on montre que $G_y$ est \emph{lisse} (6.8.1) sur $\mathbf{k}(y)$ en tous ses points~; il résultera alors de (15.6.6) que si de plus $\mathcal O_y$ est réduit, alors $f$ est lisse en tous les points de $G_y^0$ (mais pas nécessairement en tous les points de $G_y$). Ces résultats s'appliqueront par exemple dans la théorie de schémas de Picard, où nous disposerons d'un théorème général assurant que, sous certaines conditions, la fonction $z\mapsto\dim(G_z)$ est localement constante. Remarquons que c'est en vue d'applications de cette nature que les énoncés tels que (15.6.1) sont donnés pour les morphismes \emph{localement} de type fini, et non seulement pour les morphismes de type fini.

On notera aussi que dans le cas d'un $Y$-préschéma en groupes $G$, l'hypothèse $\beta$) de (15.6.6) est toujours vérifiée.
\end{remarks}

\input{ega4/ega4-15-terminal-fr.tex}

% (continued in ega4-15-terminal-fr.tex, before its first page marker)
% EGA IV-3 terminal continuation, printed pp. 242--255.
% Direct diplomatic French transcription from the NUMDAM authority.

\subsection{Appendice : Critères valuatifs de propreté locale}
\label{subsection:IV.15.7-fr}

Ce numéro donne des compléments au critère valuatif de propreté démontré dans (\textbf{II}, 7.3.10)~; il est indépendant du reste du \S~15.

\begin{definition}[15.7.1]
\label{IV.15.7.1-fr}
Soient $f:X\to Y$ un morphisme de préschémas, $y$ un point de $Y$. On dit que $f$ est propre au point $y$ s'il existe un voisinage ouvert $U$ de $y$ dans $Y$ tel que la restriction $f^{-1}(U)\to U$ de $f$ soit un morphisme propre. On dit qu'une partie $Z$ de $X$ est propre sur $Y$ au point $y$ s'il existe un voisinage ouvert $U$ de $y$ tel que $Z\cap f^{-1}(U)$ soit fermée dans $f^{-1}(U)$ et propre sur $U$ (\textbf{II}, 5.4.10) (ce qui revient à dire que pour tout sous-préschéma fermé de $X$ ayant $Z$ pour espace sous-jacent, la restriction $Z\to Y$ de $f$ à $Z$ est propre au point $y$).
\end{definition}

Soit $g:Y'\to Y$ un morphisme quelconque~; posons $X'=X\times_Y Y'$, $f'=f_{(Y')}$ et, si $p:X'\to X$ est la projection canonique, $Z'=p^{-1}(Z)$. Alors, si $Z$ est propre sur $Y$ au point $y$, $Z'$ est propre sur $Y'$ en tout point $y'$ au-dessus de $y$ (\textbf{II}, 5.4.2).

\begin{proposition}[15.7.2]
\label{IV.15.7.2-fr}
Soient $Y$ un préschéma intègre localement noethérien, $X$ un préschéma intègre, $f:X\to Y$ un morphisme séparé, dominant et de type fini, $y$ un point de $Y$. Les conditions suivantes sont équivalentes :
\begin{enumerate}[label=\emph{\alph*)}]
  \item $f$ est propre au point $y$.
  \item Si $Y_1=\operatorname{Spec}(\mathcal O_y)$, $X_1=X\times_Y Y_1$, le morphisme $f_1=f_{(Y_1)}:X_1\to Y_1$ est propre.
  \item Tout anneau de valuation discrète ayant $R(X)$ pour corps des fractions et dominant $\mathcal O$ domine un anneau local $\mathcal O_x$ de $X$ (auquel cas on a nécessairement $f(x)=y$).
\end{enumerate}
\end{proposition}

